\chapter{Obsah repozitára a reprodukcia}
\label{app:repo}

Repozitár \url{https://github.com/martinszuc/fekt-mpc-ct3-zerotrust-k8s} obsahuje všetky skripty, manifesty, namerané dáta a zdrojové texty tejto práce. Presné verzie komponentov uvádza súbor \texttt{docs/versions.md}.

\bigskip
\dirtree{%
.1 /.
.2 Makefile\DTcomment{všetky kroky experimentu ako ciele \texttt{make}}.
.2 README.md\DTcomment{požiadavky a postup reprodukcie}.
.2 infra/\DTcomment{virtual machines, k3s, Cilium, Istio, prepínanie konfigurácií}.
.2 manifests/\DTcomment{aplikácia Bookinfo a ochranné prvky po konfiguráciách}.
.2 tests/\DTcomment{skripty útokov a kontrola výsledku}.
.2 bench/\DTcomment{meranie latencie, priepustnosti a CPU}.
.2 demo/\DTcomment{scenár videozáznamu útokov}.
.2 scripts/\DTcomment{spracovanie meraní a grafy}.
.2 results/\DTcomment{výsledky útokov, dôkazy (výstupy, záznamy pcap), merania}.
.2 docs/\DTcomment{model hrozieb, plán merania, verzie, obrázky, text práce}.
}

\section*{Postup reprodukcie}

Prostredie vyžaduje počítač Mac s procesorom Apple Silicon, nástroje z Homebrew podľa súboru \texttt{README.md} a jednorazové nastavenie siete socket\_vmnet.

\begin{lstlisting}[language={},basicstyle=\ttfamily\footnotesize]
make reproduce   # cluster od nuly, útoky vo všetkých konfiguráciách
make s2a         # jedna konfigurácia: s0, s1, s2a, s2b, s2b-plain, s3
make test        # útoky proti aktuálnej konfigurácii
make bench SCENARIO="S0 S1 S2a S2b-plain S2b S3"   # meranie réžie
make plot deltas # grafy a súhrny z nameraných dát
make demo        # videozáznam útokov
make pdf         # text práce
\end{lstlisting}

Kódy konfigurácií a útokov v názvoch súborov a cieľov \texttt{make} priraďuje súbor \texttt{docs/threat-model.md}.
